% articulo-02-fuentes.tex - Motor de fuentes con LuaLaTeX
% (fontspec + luaotfload).
%
% Variables YAML (todas opcionales):
%   fuente-serif:       Arno Pro        Familia serif principal.
%   fuente-sans:        Gill Sans       Familia sans-serif.
%   fuente-mono:        Source Code Pro Familia monoespaciada.
%   fuente-mono-escala: 0.67            Escala de la monoespaciada respecto al cuerpo.
%   tipografia-sans:    (no declarar)   Si es true, el texto principal usa sans-serif.
%
% Nota sobre variantes: los nombres siguen la convención «Familia Bold/Italic/Bold Italic».
% Familias con Semibold en lugar de Bold (Caslon Pro, Garamond Pro, Jenson Pro,
% Baskerville) requieren declarar las variantes con header-includes.
%
% FALLBACK DE GLIFOS
%
%   Las fuentes OpenType complejas requieren Renderer=HarfBuzz para acceder a su
%   cobertura completa de glifos y características OpenType. Se aplica a serif y
%   sans como política general. Con HarfBuzz, sin embargo, el mecanismo
%   FallbackFonts de fontspec no funciona: HarfBuzz resuelve (o falla) la búsqueda
%   de glifos internamente, antes de que fontspec pueda interceptar y delegar.
%
%   La solución es luaotfload.add_fallback, que opera una capa por debajo: registra
%   la cadena de fallback en el motor Lua antes de que HarfBuzz procese nada, de
%   modo que HarfBuzz recibe una fuente que ya sabe delegar codepoints no cubiertos.
%
%   Cadenas declaradas:
%     seriffallback  →  Source Serif 4, Cambria Math
%     sansfallback   →  Source Sans 3, Lucida Sans Unicode
%
%   Cambria Math cubre símbolos matemáticos y letras de doble trazo (ℕ, ∈, etc.)
%   que no están en las fuentes de texto. IBM Plex Math cumple la misma
%   función para la sans.
%
%   ATENCIÓN: el nombre de la cadena (p. ej. "seriffallback") debe coincidir
%   exactamente en \directlua y en RawFeature={fallback=...}. Un error tipográfico
%   falla silenciosamente sin aviso.

\usepackage{fontspec}

% Registro de cadenas de fallback en luaotfload.
% Se declaran antes de \setmainfont para que estén disponibles en el momento
% en que fontspec construye las fuentes.
\directlua{
  luaotfload.add_fallback("seriffallback", {
    "Source Serif 4:mode=harf;",
    "Cambria Math:mode=harf;"
  })
  luaotfload.add_fallback("sansfallback", {
    "Source Sans 3:mode=harf;",
    "IBM Plex Math:mode=harf;"
  })
}

% Serif principal: Arno Pro.
% RawFeature activa la cadena arnofallback registrada arriba.
% FallbackFonts se omite: con HarfBuzz no funciona (véase nota arriba).
\setmainfont{$if(fuente-serif)$$fuente-serif$$else$Arno Pro$endif$}[
  Renderer       = HarfBuzz,
  RawFeature     = {fallback=seriffallback},
  BoldFont       = {$if(fuente-serif)$$fuente-serif$$else$Arno Pro$endif$ Bold},
  ItalicFont     = {$if(fuente-serif)$$fuente-serif$$else$Arno Pro$endif$ Italic},
  BoldItalicFont = {$if(fuente-serif)$$fuente-serif$$else$Arno Pro$endif$ Bold Italic}
]

% Sans-serif: Gill Sans.
\setsansfont{$if(fuente-sans)$$fuente-sans$$else$Gill Sans$endif$}[
  Renderer       = HarfBuzz,
  RawFeature     = {fallback=sansfallback},
  BoldFont       = {$if(fuente-sans)$$fuente-sans$$else$Gill Sans$endif$ Bold},
  ItalicFont     = {$if(fuente-sans)$$fuente-sans$$else$Gill Sans$endif$ Italic},
  BoldItalicFont = {$if(fuente-sans)$$fuente-sans$$else$Gill Sans$endif$ Bold Italic}
]

% Monoespaciada: Source Code Pro.
% Sin fallback: los caracteres especiales no son esperables en código.
\setmonofont{$if(fuente-mono)$$fuente-mono$$else$Source Code Pro$endif$}[
  BoldFont       = {$if(fuente-mono)$$fuente-mono$$else$Source Code Pro$endif$ Bold},
  ItalicFont     = {$if(fuente-mono)$$fuente-mono$$else$Source Code Pro$endif$ Italic},
  BoldItalicFont = {$if(fuente-mono)$$fuente-mono$$else$Source Code Pro$endif$ Bold Italic},
  Scale          = $if(fuente-mono-escala)$$fuente-mono-escala$$else$0.67$endif$
]

% Selector de familia principal desde metadato YAML.
% tipografia-sans: true  →  texto en sans-serif
% (sin declarar)         →  texto en serif (por defecto)
$if(tipografia-sans)$
\renewcommand{\familydefault}{\sfdefault}
$endif$
